\chapter{Intra-articular Steroid Injection}

\section*{Overview}
Intra-articular steroid injection is the delivery of a corticosteroid directly into a joint space to reduce local inflammation and pain. It is commonly used for inflammatory and degenerative joint diseases. The exact approach, setting, and preparation depend on the indication, anatomy, and the patient's health.

\section*{Alternative Names}
Joint injection; corticosteroid joint injection; intra-articular corticosteroid; IA steroid injection

\section*{Principle}
Corticosteroids suppress local inflammation by inhibiting leukocyte migration, cytokine release, and prostaglandin synthesis within the joint. Delivering the drug directly into the joint achieves a high local concentration with limited systemic exposure, and the effect typically lasts weeks to months.

\section*{History}
Intra-articular corticosteroid injection was introduced by Hollander and colleagues in the early 1950s and remains a cornerstone of local therapy for inflammatory arthritis.

\section*{Why It Is Done}
Ordered for symptomatic relief of osteoarthritis, rheumatoid arthritis flare, gout, pseudogout, and other inflammatory arthropathies, and to bridge care while systemic therapy takes effect.

\section*{Is This a Primary or Secondary Test or Procedure}
A secondary, adjunctive treatment for joint inflammation; it provides symptom relief rather than disease modification.

\section*{How It Is Performed}
The joint is identified by palpation or ultrasound guidance, and the skin is cleaned with antiseptic. A needle is inserted through the skin into the joint space, synovial fluid may be aspirated for analysis, and a corticosteroid, sometimes with local anesthetic, is injected before the needle is withdrawn.

\section*{Preparation}
Informed consent, review of anticoagulation and bleeding risk, and confirmation that the joint is not septic are needed. The patient should report signs of active infection, and steroid should generally be held for one to two weeks before elective surgery on a joint.

\section*{Contraindications and Precautions}
Septic arthritis, bacteremia, skin infection over the joint, an unstable or prosthetic joint, and allergy to the injected agents are contraindications. Precautions include limiting injection frequency, usually no more than every three months and no more than three to four times per joint per year, and exercising caution in diabetic patients.

\section*{Risks}
Risks include post-injection flare, local infection, bleeding or bruising, tendon rupture, skin or soft tissue atrophy, transient glucose elevation, and rarely intra-articular crystal deposition with cartilage damage from repeated injections.

\section*{Aftercare and Recovery}
The joint is rested for 24 to 48 hours and iced, and the patient is advised that a brief post-injection flare may occur. Relief typically develops within a few days and may last weeks to several months.

\section*{Outcome and Follow-up}
Durability of pain relief and improved joint mobility indicate a response, while early recurrence suggests more severe disease and the need for systemic or surgical therapy. A hot, red, or swollen joint after injection warrants evaluation for infection.

\section*{Expected Results}
Marked reduction in joint pain, swelling, and stiffness with improved range of motion for weeks to months after injection indicates a good result.

\section*{Follow-up}
Physical therapy and strengthening are continued, and systemic disease-modifying therapy is optimized for inflammatory arthritis. Repeat injection or arthroplasty is considered when relief is short-lived or disease progresses.

\section*{When to Seek Medical Care}
Follow the procedure team's specific instructions. Seek urgent medical assessment for severe or worsening pain, uncontrolled bleeding, fever or signs of infection, trouble breathing, chest pain, fainting, new weakness or confusion, inability to urinate, or any rapidly worsening symptom. The appropriate warning signs depend on the procedure and the patient's medical condition.

\section*{References}
\begin{enumerate}
  \item MedlinePlus. Steroid injection (intra-articular). U.S. National Library of Medicine; 2025.
  \item Current Medical Diagnosis and Treatment. McGraw-Hill; 2025.
\end{enumerate}

\clearpage
